% Options for packages loaded elsewhere
\PassOptionsToPackage{unicode}{hyperref}
\PassOptionsToPackage{hyphens}{url}
\PassOptionsToPackage{space}{xeCJK}
%
\documentclass[
]{ctexart}
\usepackage{amsmath,amssymb}
\usepackage{iftex}
\ifPDFTeX
  \usepackage[T1]{fontenc}
  \usepackage[utf8]{inputenc}
  \usepackage{textcomp} % provide euro and other symbols
\else % if luatex or xetex
  \usepackage{unicode-math} % this also loads fontspec
  \defaultfontfeatures{Scale=MatchLowercase}
  \defaultfontfeatures[\rmfamily]{Ligatures=TeX,Scale=1}
\fi
\usepackage{lmodern}
\ifPDFTeX\else
  % xetex/luatex font selection
  \ifXeTeX
    \usepackage{xeCJK}
    \setCJKmainfont[]{Songti SC}
          \fi
  \ifLuaTeX
    \usepackage[]{luatexja-fontspec}
    \setmainjfont[]{Songti SC}
  \fi
\fi
% Use upquote if available, for straight quotes in verbatim environments
\IfFileExists{upquote.sty}{\usepackage{upquote}}{}
\IfFileExists{microtype.sty}{% use microtype if available
  \usepackage[]{microtype}
  \UseMicrotypeSet[protrusion]{basicmath} % disable protrusion for tt fonts
}{}
\makeatletter
\@ifundefined{KOMAClassName}{% if non-KOMA class
  \IfFileExists{parskip.sty}{%
    \usepackage{parskip}
  }{% else
    \setlength{\parindent}{0pt}
    \setlength{\parskip}{6pt plus 2pt minus 1pt}}
}{% if KOMA class
  \KOMAoptions{parskip=half}}
\makeatother
\usepackage{xcolor}
\usepackage[margin=1in]{geometry}
\setlength{\emergencystretch}{3em} % prevent overfull lines
\providecommand{\tightlist}{%
  \setlength{\itemsep}{0pt}\setlength{\parskip}{0pt}}
\setcounter{secnumdepth}{-\maxdimen} % remove section numbering
\def\ba{\begin{aligned}}
\def\ea{\end{aligned}}
\def\bm{\begin{bmatrix}}
\def\em{\end{bmatrix}}
\def\bc{\begin{cases}}
\def\ec{\end{cases}}

\def\ve{\varepsilon}
\def\s{\sigma}
\newcommand{\0}{\textbf 0}
\newcommand{\1}{\textbf 1}
\def\I{\textbf I}
\def\X{\textbf X}
\def\x{\textbf x}
\def\y{\textbf y}
\def\b{\textbf b}
\def\w{\textbf w}

\def\plim{\text{plim }}
\providecommand{\var}{}\renewcommand{\var}[2][]{\text{Var}_{#1}\br{#2}}
\providecommand{\e}{}\renewcommand{\e}[2][]{\text{E}_{#1}\br{#2}}
\providecommand{\cov}{}\renewcommand{\cov}[3][]{\text{Cov}_{#1}\br{#2,#3}}
\def\se#1{\text{se}\pr{#1}}
\def\estvar#1{\text{Est.Var}\br{#1}}
\def\estcov#1#2{\text{Est.Cov}\br{#1,#2}}
\def\asyvar#1{\text{Asy.Var}\br{#1}}
\def\asycov#1#2{\text{Asy.Cov}\br{#1,#2}}
\def\estasyvar#1{\text{Est.Asy.Var}\br{#1}}
\def\too{\Longrightarrow}
\def\sq{\sqrt}
\def\ss{\frac{e'e}{n-K}}
\providecommand{\dev}{}\renewcommand{\dev}[2][i]{({#2}_{#1}-\bar{#2})}
\def\inv#1{{#1}^{-1}}
\providecommand{\add}{}\renewcommand{\add}[1][i]{\sum_{#1 = 1}^n}
\providecommand{\ply}{}\renewcommand{\ply}[1][i]{\prod_{#1 = 1}^n}
\providecommand{\l}{}\renewcommand{\l}[1][(]{\left#1}
\providecommand{\r}{}\renewcommand{\r}[1][)]{\right#1}
\def\br#1{\left[#1\right]}
\def\pr#1{\left(#1\right)}
\def\brace#1{\left\{#1\right\}}
\providecommand{\seq}{}\renewcommand{\seq}[1][x]{{#1}_1,{#1}_2,\cdots,{#1}_{n}}
\providecommand{\seqadd}{}\renewcommand{\seqadd}[1][x]{{#1}_1+{#1}_2+\cdots+{#1}_{n}}
\def\dto{\overset{\text d}{\to}}
\def\asim{\overset{\text a}{\sim}}

\providecommand{\P}{}\renewcommand{\P}[1][X]{{#1}({#1}'{#1})^{-1}{#1}'}
\providecommand{\M}{}\renewcommand{\M}[1][X]{\I-{#1}({#1}'{#1})^{-1}{#1}'}
\providecommand{\bls}{}\renewcommand{\bls}[1][X]{(#1'#1)^{-1}#1'y}
\providecommand{\bx}{}\renewcommand{\bx}[1][X]{(#1'#1)^{-1}#1'}
\providecommand{\biv}{}\renewcommand{\biv}[1][Z]{(#1'X)^{-1}#1'y}
\providecommand{\XX}{}\renewcommand{\XX}[1][X]{(#1'#1)^{-1}}
\def\Q{\frac{X'X}{n}}
\providecommand{\Xe}{}\renewcommand{\Xe}[1][X]{\frac{{#1}'\ve}{n}}
\providecommand{\RX}{}\renewcommand{\RX}[1][X]{R({#1}'{#1})^{-1}R'}
\providecommand{\m}{}\renewcommand{\m}[1][b]{R{#1}-q}
\def\O{X'\Omega X}

\providecommand{\colvec}{}\renewcommand{\colvec}[2][n]{
\begin{bmatrix}
{#2}_1\\
{#2}_2\\
\vdots \\
{#2}_{#1}
\end{bmatrix}
}

\providecommand{\mx}{}\renewcommand{\mx}[3][a]{
\begin{bmatrix}
a_{11} & a_{12} & \cdots & a_{1#3}\\
a_{21} & a_{22} & \cdots & a_{2#3}\\
\vdots & \vdots & \ddots & \vdots \\
a_{{#2}1} & a_{{#2}2} & \cdots & a_{{#2}{#3}}
\end{bmatrix}
}
\providecommand{\diag}{}\renewcommand{\diag}[2][n]{
\begin{bmatrix}
{#2}_{1} & 0 & \cdots & 0\\
0 & {#2}_{2} & \cdots & 0\\
\vdots & \vdots & \ddots & \vdots \\
0 & 0 & \cdots & {#2}_{{#1}}
\end{bmatrix}
}


\ifLuaTeX
  \usepackage{selnolig}  % disable illegal ligatures
\fi
\IfFileExists{bookmark.sty}{\usepackage{bookmark}}{\usepackage{hyperref}}
\IfFileExists{xurl.sty}{\usepackage{xurl}}{} % add URL line breaks if available
\urlstyle{same}
\hypersetup{
  hidelinks,
  pdfcreator={LaTeX via pandoc}}

\author{}
\date{}

\begin{document}

\renewcommand*\contentsname{目录}
{
\setcounter{tocdepth}{3}
\tableofcontents
}
\hypertarget{causal-relationship}{%
\section{Causal Relationship}\label{causal-relationship}}

\hypertarget{dag}{%
\subsection{DAG}\label{dag}}

\begin{quote}
\begin{itemize}
\tightlist
\item
  A graph that is both \textbf{directed} and \textbf{acylic} is called a
  \textbf{directed acyclic graph} (DAG).
\item
  A back-door path from node \(A\) to node \(B\) is a path from \(A\) to
  \(B\) that starts with an incoming arrow into \(A\) and ends with an
  incoming arrow into \(B\).
\end{itemize}
\end{quote}

Basic patterns of causal relationships:

\begin{enumerate}
\def\labelenumi{\arabic{enumi}.}
\tightlist
\item
  Chains (causal pathway)
\item
  Forks (confounding pathway)
\item
  Inverted forks (colliding pathway)
\end{enumerate}

\hypertarget{strategies-to-estimate-causal-effects}{%
\subsection{Strategies to estimate causal
effects}\label{strategies-to-estimate-causal-effects}}

\begin{enumerate}
\def\labelenumi{\arabic{enumi}.}
\tightlist
\item
  condition on variables (that block all back-door paths from the causal
  variable to the outcome variable).

  \begin{itemize}
  \tightlist
  \item
    However, conditioning on a collider variable does not simplify the
    original graph but rather adds complications by creating new
    association.
  \end{itemize}
\item
  use exogenous variation in an appropriate instrumental variable (to
  isolate covariation in the causal and outcome variables).
\item
  establish an isolated and exhaustive mechanism (that relates the
  causal variable to the outcome variable).
\end{enumerate}

\hypertarget{back-door-criterion}{%
\subsection{Back-door criterion}\label{back-door-criterion}}

\begin{quote}
The causal effect is identified by conditioning on a set of variables in
\(Z\) if and only if all back-door paths between the causal variable and
the outcome variable are blocked after conditioning on \(Z\).

All back-door paths are blocked by \(Z\) (i.e., \(d\)-separated by
\(Z\)) if and only if each back-door path:

\begin{enumerate}
\def\labelenumi{\arabic{enumi}.}
\tightlist
\item
  contains a chain of mediation \(A\to C\to B\), where the middle
  variable \(C\) is in \(Z\);
\item
  contains a fork of mutual dependence \(A\leftarrow C\to B\), where the
  middle variable \(C\) is in \(Z\);
\item
  contains an inverted fork of mutual causation \(A\to C\leftarrow B\),
  where the middle variable \(C\) and all of \(C\)'s descendants are not
  in \(Z\).
\end{enumerate}
\end{quote}

Back-door path: A path between any causally ordered sequence of two
variables that includes a directed edge \(\to\) that points to the first
variable.

\hypertarget{front-door-criterion}{%
\subsection{Front-door criterion}\label{front-door-criterion}}

\begin{quote}
One can consistently estimate the effect of a causal variable on an
outcome variable by estimating the effect as it propagates through an
isolated and exhaustive mechanism.

\begin{itemize}
\tightlist
\item
  Isolated: none of the components of unblockable back-door paths have
  direct effects on the mechanism.
\item
  Exhaustive: All front-door paths are identifiable.
\end{itemize}
\end{quote}

Prove front-door criterion from back-door criterion:

\begin{itemize}
\tightlist
\item
  \(D\to M\): There are no back-door paths between \(D\) and \(M\).
\item
  \(M\to Y\): There are some back-door paths via \(D\), however all of
  them can be blocked by \(D\).
\end{itemize}

\hypertarget{causal-inference}{%
\section{Causal Inference}\label{causal-inference}}

Average treatment effect (ATE) \[
\tau_{ATE} = \e{Y_{i}\pr 1-Y_{i}\pr 0}
\] Average treatment effect on the treated (ATET) \[
\tau_{ATET} = \e{Y_{i}\pr 1-Y_{i}\pr 0|D_i = 1}
\] Average treatment effect on the untreated (ATENT) \[
\tau_{ATENT} = \e{Y_{i}\pr 1-Y_{i}\pr 0|D_i = 0}
\]

\begin{quote}
ATE is a \textbf{weighted average} of the ATET and the ATENT \[
\tau_{ATE} = p\pr{D_i = 1}\cdot \tau_{ATET}+p\pr{D_i = 0}\cdot \tau_{ATENT}
\]
\end{quote}

\begin{quote}
Conditional Independence Assumption: \(Y_{\pr{d}}\perp D_i |X_i\).

Conditional Mean Independence:
\(\e{Y_i|x,D} = \e{Y_i|x},\e{Y_0|x,D} = \e{Y_0|x}\).
\end{quote}

\hypertarget{did}{%
\section{DID}\label{did}}

Idea of DID: Compare the outcomes of the units that are affected by the
policy change and those who are not affected before and after the policy
was enacted.

\hypertarget{did-with-various-forms}{%
\subsection{DID with various forms}\label{did-with-various-forms}}

\begin{quote}
\begin{enumerate}
\def\labelenumi{\arabic{enumi}.}
\item
  Different treatment time \[
  y_{it} = \alpha_0+\beta D_{it}+\mu_i+f_t+\ve_{it}
  \]
\item
  Different treatment magnitude \[
  y_{it} = \alpha_0+\beta \pr{x_i\cdot d}+\mu_i+f_t+\ve_{it}
  \]
\item
  A general form \[
  y_{it} = \alpha_0 + \beta_t \pr{x_i\cdot d_{it}}+\mu_i+f_t+\ve_{it}
  \]
\end{enumerate}
\end{quote}

\hypertarget{did-with-matching}{%
\subsection{DID with matching}\label{did-with-matching}}

\begin{itemize}
\tightlist
\item
  It can be seen as a non-parametric DID, reweighting observations
  according to a weighting function dependent on the specific matching
  approach adopted.
\item
  Advantage: it does not require the imposition of the
  linear-in-parameters form of the outcome equation.
\end{itemize}

For panel data \[
\widehat{ATET}_{\mathit{M-DID}} = \frac1{N_1}\sum_{i\in T}\br{\pr{Y_{i1}^T-Y_{i0^T}-\sum_{j\in C\pr i}h\pr{i,j}\pr{Y_{j1}^C-y_{j0}^C}}}
\]

\hypertarget{parallel-trend-assumption}{%
\subsection{Parallel trend assumption}\label{parallel-trend-assumption}}

\begin{itemize}
\tightlist
\item
  Test the assumption for periods before the treatment.
\item
  Allows for heterogeneity in time trend.
\end{itemize}

\hypertarget{iv}{%
\section{IV}\label{iv}}

\hypertarget{problems-of-iv-estimation}{%
\subsection{Problems of IV estimation}\label{problems-of-iv-estimation}}

\begin{enumerate}
\def\labelenumi{\arabic{enumi}.}
\tightlist
\item
  Inconsistency of IV: What if the instrument is not fully exogenous for
  the outcome. The bias increases as the covariance of \(Z\) and \(X\)
  decreases.
\item
  Lower efficiency of IV: The variance of IV is always greater than that
  of OLS.
\item
  Small-sample bias: The bias of 2SLS when one cannot invoke the usual
  asymptotic results.
\end{enumerate}

\hypertarget{estimation-with-a-binary-iv}{%
\subsection{Estimation with a binary
IV}\label{estimation-with-a-binary-iv}}

\[
\hat \beta = \frac{\e{y|z = 1}-\e{y|z= 0}}{\e{d|z = 1}-\e{d|z = 0}}
\]

\begin{itemize}
\tightlist
\item
  An IV estimator is a ratio that is a joint projection of \(y\) and
  \(d\) onto a third dimension of \(z\).
\item
  An IV estimator isolates a specific portion of the covariation in
  \(D\) and \(Y\).

  \begin{itemize}
  \tightlist
  \item
    By using the fitted values of the endogenous regressor from the
    first-stage regression, our regression now uses only the exogenous
    variation in the regressor due to the instrumental variable itself.
  \item
    Instrumental variable therefore reduces the variation in the data,
    but that variation is exogenous.
  \end{itemize}
\item
  A binary IV identifies only the average causal effect for compliers.
  (the subset of all students who would attend a private school if given
  a voucher but who would not attend a private school in the absence of
  a voucher)
\end{itemize}

\hypertarget{late}{%
\subsection{LATE}\label{late}}

The LATE theorem:

\begin{itemize}
\tightlist
\item
  Independece
\item
  Exclusion
\item
  First stage
\item
  Monotonicity
\end{itemize}

LATE is not the same as ATT: \[
\ba
\underbrace{\e{Y_i^1-Y_i^0|D_i = 1}}_{ATT} = & \underbrace{\e{Y_i^1-Y_i^0|D_i = 1}\cdot p\pr{D_i^0 = 1|D_i = 1}}_{\text{Effect on Always-Takers}}\\
& + \underbrace{\e{Y_i^1-Y_i^0|D_i = 1}\cdot p\pr{D_i^1>D_i^0,Z_i = 1|D_i = 1}}_{\text{Effect on Compliers}}
\ea
\] ATT is a weighted average of the effects on always-takers and
compliers.

\hypertarget{psm}{%
\section{PSM}\label{psm}}

\hypertarget{balancing-scores}{%
\subsection{Balancing scores}\label{balancing-scores}}

\begin{quote}
A balancing score \(b\pr x\) is a function of the covariates such that
\[
W_i\perp X_i|b\pr{X_i}
\]
\end{quote}

The idea is to find lower-dimensional functions of the covariates that
suffice for the removing the bias associated with differences in the
pre-treatment variables.

Property: If assignment to treatment is unconfounded given the full set
of covariates, then assignment is also unconfounded conditioning only on
a balancing score.

\begin{itemize}
\tightlist
\item
  The propensity score provides the biggest benefit in terms of reducing
  the number of variables we need to adjust for.
\item
  Units stratified according to the propensity score should be
  indistinguishable in terms of their \(X\).
\end{itemize}

\hypertarget{identification-assumptions-of-matching}{%
\subsection{Identification assumptions of
matching}\label{identification-assumptions-of-matching}}

\begin{itemize}
\tightlist
\item
  Conditional independence (unconfoundedness): There exists a set of
  observable covariates such that after controlling for these
  covariates, treatment assignment is independent of potential outcomes.
\item
  Common support: For each value of \(X\), there is a positive
  probability of being both treated and untreated.
\end{itemize}

\hypertarget{decisions-to-make-in-matching}{%
\subsection{Decisions to make in
matching}\label{decisions-to-make-in-matching}}

\begin{itemize}
\tightlist
\item
  Variables to be used for matching.
\item
  Distances for measuring closeness in covariates.
\item
  The number of controls.
\item
  Caliper.
\item
  Stratification: For covariates that should be matched exactly, forming
  the same strata depending on such covariates in the treatment and
  control groups.
\item
  Greedy/Nongreedy. (Greedy matching: A control is matched to only one
  treated at most.)
\end{itemize}

\hypertarget{basic-matching-estimators}{%
\subsection{Basic matching estimators}\label{basic-matching-estimators}}

\[
\hat\tau = \frac1{N_1}\sum_{t\in T}\pr{Y_t-\sum_{c\in\pr{C}}w_{ct}Y_c}
\]

However, it is hard to find the asymptotic distribution of a matching
estimator, because selecting a comparison group \(C_t\) for treated
\(i\) involves all observations. This implies dependence across the
paired differences to make the iid presumption is false. Ignoring this
dependence is likely to underestimate the asymptotic variance.

\hypertarget{assessing-overlap-in-covariate-distribution}{%
\subsection{Assessing overlap in covariate
distribution}\label{assessing-overlap-in-covariate-distribution}}

It is useful to assess the degree of overlap in the covariate
distributions.

\begin{itemize}
\tightlist
\item
  The difference in locations: normalized difference
\item
  The difference in dispersions: the logarithm of the ratio of standard
  deviations
\item
  Direct measures of overlap: investigate the fraction of the treated
  units have covariates values that are in the tails of the distribution
  of the covariate values for the controls.
\end{itemize}

\hypertarget{rct}{%
\section{RCT}\label{rct}}

\hypertarget{definition-of-treatment-effects}{%
\subsection{Definition of treatment
effects}\label{definition-of-treatment-effects}}

Individual treatment effect is the difference between the potential
outcomes of it being treated minus those when it not being treated.

There are two important aspects of the definition of a causal effect:

\begin{itemize}
\tightlist
\item
  The definition of the causal effect depends on the \textbf{potential
  outcomes}, but it does not depend on which outcome is actually
  observed.
\item
  The causal effect is the comparison of potential outcomes, for the
  same unit, at the same moment after the treatment.
\end{itemize}

In experimental and quasi-experimental designs, the treatment effect is
generally estimated by the \textbf{counterfactual approach}.
Counterfactual causality takes the form of a comparison between the
outcome of a unit when the unit is treated in a certain way, and the
outcome of the same unit when it is not treated.

\hypertarget{bias-in-estimation-of-treatment-effects}{%
\subsection{Bias in estimation of treatment
effects}\label{bias-in-estimation-of-treatment-effects}}

We can use observed outcomes to estimate the treatment effect for unit
\(i\): \[
Y_i - Y_j = Y_i(1) - Y_j(0) = \underbrace{Y_i(1) - Y_i(0)}_{\gamma_i} + \underbrace{Y_i(0) - Y_j(0)}_{\text{Selection Bias}}
\] Bias in estimation of average treatment effect using observed
outcomes:

\begin{itemize}
\item
  Bias in estimation of ATT using ``naive'' estimator: \[
  T_1 - C_0 = \underbrace{T_1-T_0}_{ATT}+\underbrace{T_0-C_0}_\text{bias}
  \]
\item
  Bias in estimation of ATU using ``naive'' estimator: \[
  T_1 - C_0 = \underbrace{C_1-C_0}_{ATU}+\underbrace{T_1-C_1}_\text{bias}
  \]
\item
  Bias in estimation of ATE using ``naive'' estimator: \[
  T_1 - C_0 = \underbrace{\omega (T_1-T_0) + (1-\omega)(C_1-C_0)}_{ATE}+\underbrace{\omega (T_0 -C_0) + (1-\omega)(T_1-C_1)}_\text{bias}
  \]
\end{itemize}

\hypertarget{properties-of-assignment-mechanism}{%
\subsection{Properties of assignment
mechanism}\label{properties-of-assignment-mechanism}}

\begin{enumerate}
\def\labelenumi{\arabic{enumi}.}
\tightlist
\item
  \textbf{Individualistic} assignment: It limits the dependence of the
  treatment assignment for unit \(i\) on the outcomes and assignments
  for other units.
\item
  \textbf{Probabilistic} assignment: Every unit has positive probability
  of being assigned to the treatment or the control group.
\item
  \textbf{Unconfounded} assignment: The assignment mechanism does not
  depend on the potential outcomes.
\end{enumerate}

Under these assumptions, we can use the naive estimator to estimate
treatment effects.

\hypertarget{rdd}{%
\section{RDD}\label{rdd}}

\hypertarget{identification-at-cutoff}{%
\subsection{Identification at cutoff}\label{identification-at-cutoff}}

Formally, RD refers to a regression model \[
\e{Y|S} = \beta_d\e{D|S}+m\pr{S}
\] \hspace{0pt} where \(\e{D|S}\) is discontinuous at \(0\), and
\(m\pr{S}\) is an unknown function continuous at \(0\).

\[
\beta_d = \frac{\e{Y|0^+}-\e{Y|0^-}}{\e{D|0^+}-\e{D|0^-}}
\] The RD identification condition has two components:

\begin{itemize}
\tightlist
\item
  The discontinuity of \(\e{D|S}\): \(\e{D|S}\) has a break at \(0\),
  that is, \(\e{D|0^+}-\e{D|0^-} \ne 0\).
\item
  The continuity of \(m\pr{S}\) at \(0\):
  \(m\pr{S} = \e{Y|S}-\beta_d\e{D|S}\) is continuous at \(0\).
\end{itemize}

\hypertarget{main-features-of-rd}{%
\subsection{Main features of RD}\label{main-features-of-rd}}

In RD, the discontinuity of \(\e{D|S}\) at \(0\) comes typically from an
institutional feature and thus it is visible, but the continuity of
\(m\pr{S}\) is an assumption.

This continuity leads to three main features of RD:

\begin{itemize}
\tightlist
\item
  \textbf{Local randomization}.

  \begin{itemize}
  \tightlist
  \item
    In SRD with \(D = 1\br{S\ge 0}\)
  \end{itemize}
\item
  \textbf{Robustness of RD to the endogeneity of \(D\) through \(S\).}

  \begin{itemize}
  \tightlist
  \item
    Data generating process: \(Y_i = \beta_d D_i+m\pr{S_i}+U_i\), where
    \(\e{U|S} = 0\).
  \item
    RD allows \(\e{U|S}\) to be a nontrivial function of \(S\), as long
    as \(\e{U|S}\) is continuous at \(0\); such a \(\e{U|S}\) can be
    merged into \(m\pr{S}\).
  \item
    The capacity to allow \(\e{U|S}\) to be a function of \(S\) makes RD
    estimators robust to the endogeneity of \(D\) through \(S\).
  \item
    \(D\) in FRD affected by both \(S\) and \(\ve\) can be still related
    to \(U\) through \(\ve\), and RD estimators are not robust to the
    \(D\) endogeneity through \(\ve\).
  \end{itemize}
\item
  \textbf{Ignorability of covariates other than \(S\).}

  \begin{itemize}
  \tightlist
  \item
    As long as the covariates are continuous at the cutoff point, the
    covariates do not affect the procedure of identification.
  \end{itemize}
\end{itemize}

\hypertarget{data-requirement-of-rdd}{%
\subsection{Data requirement of RDD}\label{data-requirement-of-rdd}}

\begin{enumerate}
\def\labelenumi{\arabic{enumi}.}
\tightlist
\item
  Three basic variables:

  \begin{itemize}
  \tightlist
  \item
    Assignment variable (or called forcing variable, running variable):
    \(S\)
  \item
    Cutoff point: \(C\)
  \item
    Outcome variable: \(Y\)
  \end{itemize}
\item
  Values of the assignment variable cannot be manipulated around the
  cutoff point (to reach local randomization).
\item
  Cutoff point is not affected by the assignment variable (one cannot
  choose a cutoff point to change the assignment).
\item
  Individual characteristics are unaffected around the cutoff point,
  other than the assignment status.
\end{enumerate}

\hypertarget{specification-tests-of-rdd}{%
\subsection{Specification tests of
RDD}\label{specification-tests-of-rdd}}

\begin{enumerate}
\def\labelenumi{\arabic{enumi}.}
\tightlist
\item
  The breaks of \(\e{Y|S}\) and \(\e{D|S}\).
\item
  The continuity of \(m\pr S\).
\end{enumerate}

\hypertarget{challenges-to-rdd}{%
\subsection{Challenges to RDD}\label{challenges-to-rdd}}

\begin{itemize}
\tightlist
\item
  Treatment is not as good as randomly assigned around the cutoff when
  agents are able to manipulate their running variable scores.
\item
  Some other unobservable characteristic changes at the threshold, and
  this has a direct effect on the outcome.
\end{itemize}

\end{document}
